\section{Conclusion}
\label{sec:conclusion}

A model holding a secret leaks it into the content decisions it makes, whenever it makes them.
Handing \gemma a plan helps only when the plan was made without the secret: leakage of a secret word into the prose then drops from 84.1\% to about chance, while the secret remains fully decodable inside the model.
A plan the model writes for itself absorbs the leak and passes it on.
Context that fixes no content changes little, and the leak needs the secret to stay in context while the body of the story is written.
For plot twists the direction is the same and the evidence weaker.

Several follow-ups would strengthen these findings.
The saved stories should be re-judged with frontier guessers and with humans, especially in the blind-outline and twist conditions.
Masking attention from story positions to the secret tokens, layer by layer, would locate the reads that the swap experiment says are necessary.
The twist experiment should be run on whole stories with a larger writer.
Finally, our results suggest a two-context recipe that we did not test end to end: plan in a context without the secret, then write in one with it.
